\chapter{Logique, ensembles et applications}\label{ch:b1:logic}

Jusqu'ici, les démonstrations ont été menées avec une idée informelle
mais honnête de ce que «~démontrer~» veut dire. Ce premier chapitre de
mathématiques du supérieur explicite les règles du jeu~: ce qu'est une
\omterm{def:b1:logic:statement}{assertion} mathématique, comment les connecteurs et les quantificateurs
combinent les \omterm{def:b1:logic:statement}{assertions}, quels coups sont autorisés dans une
démonstration --- puis construit, sur cette base, les deux langages
universels des mathématiques~: les \omterm{def:b1:logic:sets}{ensembles} et les \omterm{def:b1:logic:map}{applications}.

\section{Assertions et connecteurs}

\begin{definition}[Assertion, connecteurs]\label{def:b1:logic:statement}
Une \emph{assertion}\index{assertion} (ou \emph{proposition}) est une
phrase qui est soit vraie (V), soit fausse (F) --- exactement l'une des
deux. À partir de deux assertions $P$ et $Q$, on forme~:
\begin{itemize}
  \item la \emph{négation} $\lnot P$ («~non $P$~»), vraie exactement
        quand $P$ est fausse~;
  \item la \emph{conjonction} $P \land Q$ («~$P$ et $Q$~»), vraie
        exactement quand les deux sont vraies~;
  \item la \emph{disjonction} $P \lor Q$ («~$P$ ou $Q$~»), vraie
        exactement quand l'une au moins est vraie (ce «~ou~» est
        inclusif)~;
  \item l'\emph{implication} $P \implies Q$, fausse exactement quand $P$
        est vraie et $Q$ fausse~;
  \item l'\emph{équivalence} $P \iff Q$, vraie exactement quand $P$ et
        $Q$ ont la même valeur de vérité.
\end{itemize}
\end{definition}

\begin{remark}
La table de vérité de $P \implies Q$ mérite qu'on s'y arrête~: quand $P$
est fausse, $P \implies Q$ est \emph{vraie}, quel que soit $Q$. «~Si
$2 < 1$ alors $0 = 5$~» est une implication vraie. Une implication
n'affirme rien sur ce qui se passe quand son hypothèse est en défaut.
\end{remark}

\begin{proposition}[Règles de calcul sur les assertions]\label{prop:b1:logic:rules}
Pour toutes \omterm{def:b1:logic:statement}{assertions} $P$, $Q$, $R$~:
\begin{enumerate}
  \item $\lnot(\lnot P) \iff P$~;
  \item \omterm{prop:b1:logic:rules}{lois de De Morgan}\index{lois de De Morgan}~:
        $\lnot(P \land Q) \iff (\lnot P) \lor (\lnot Q)$ et
        $\lnot(P \lor Q) \iff (\lnot P) \land (\lnot Q)$~;
  \item $(P \implies Q) \iff \bigl((\lnot P) \lor Q\bigr)$, d'où
        $\lnot(P \implies Q) \iff P \land (\lnot Q)$~;
  \item contraposition\index{contraposition}~:
        $(P \implies Q) \iff \bigl((\lnot Q) \implies (\lnot P)\bigr)$~;
  \item $(P \iff Q) \iff \bigl((P \implies Q) \land (Q \implies
        P)\bigr)$~;
  \item distributivité~: $P \land (Q \lor R) \iff (P \land Q) \lor
        (P \land R)$ et $P \lor (Q \land R) \iff (P \lor Q) \land
        (P \lor R)$.
\end{enumerate}
\end{proposition}

\begin{proof}
Chaque équivalence se vérifie en comparant des tables de vérité~: deux
\omterm{def:b1:logic:statement}{assertions} composées construites à partir de $P$, $Q$, $R$ sont
équivalentes exactement lorsqu'elles prennent la même valeur de vérité
dans chacun des (quatre ou huit) cas. Écrivons une table en entier,
pour la première loi de De Morgan~:
\begin{center}
\begin{tabular}{cc|c|c|cc|c}
$P$ & $Q$ & $P \land Q$ & $\lnot(P \land Q)$ & $\lnot P$ & $\lnot Q$ &
$(\lnot P) \lor (\lnot Q)$ \\
\hline
V & V & V & F & F & F & F \\
V & F & F & V & F & V & V \\
F & V & F & V & V & F & V \\
F & F & F & V & V & V & V \\
\end{tabular}
\end{center}
Les colonnes $4$ et $7$ coïncident, ce qui démontre la loi. Pour la
contraposition, un raccourci verbal va plus vite~: $P \implies Q$ est
fausse exactement dans le cas ($P$ vraie, $Q$ fausse), et
$(\lnot Q) \implies (\lnot P)$ est fausse exactement dans le cas
($\lnot Q$ vraie, $\lnot P$ fausse), c'est-à-dire ($Q$ fausse, $P$
vraie) --- le même cas unique, donc les deux implications ont des tables
identiques. Les autres règles se vérifient de la même façon~; notons que
(3) ramène toute implication à une disjonction, si bien que (2) produit
mécaniquement la règle de négation $\lnot(P \implies Q) \iff P
\land (\lnot Q)$~: pour contredire une implication, il faut exhiber un
cas où l'hypothèse est vérifiée et la conclusion en défaut.
\end{proof}

\section{Quantificateurs}

\begin{definition}[Quantificateurs]\label{def:b1:logic:quantifiers}
Soit $P(x)$ une propriété d'un élément $x$ d'un \omterm{def:b1:logic:sets}{ensemble} $E$.
\begin{itemize}
  \item $\forall x \in E,\ P(x)$ («~pour tout $x$ de $E$, $P(x)$~») est
        vraie quand tout élément de $E$ vérifie $P$~;
  \item $\exists x \in E,\ P(x)$ («~il existe $x$ dans $E$ tel que
        $P(x)$~») est vraie quand au moins un élément de $E$ vérifie
        $P$.
\end{itemize}
On écrit $\exists!$ pour «~il existe un unique~».
\end{definition}

\begin{proposition}[Négation des quantificateurs]\label{prop:b1:logic:negquant}
\[
\lnot\bigl(\forall x \in E,\ P(x)\bigr) \iff
\exists x \in E,\ \lnot P(x),
\qquad
\lnot\bigl(\exists x \in E,\ P(x)\bigr) \iff
\forall x \in E,\ \lnot P(x).
\]
\end{proposition}

\begin{proof}
Démontrons la première équivalence dans les deux sens~; la seconde
est symétrique. Si $\forall x \in E,\ P(x)$ est fausse, alors les
éléments de $E$ ne vérifient pas tous $P$~: l'\omterm{def:b1:logic:sets}{ensemble}
$A = \{x \in E : \lnot P(x)\}$ ne peut être vide, et n'importe lequel
de ses éléments témoigne de
$\exists x \in E,\ \lnot P(x)$. Réciproquement, si un certain $x_0 \in
E$ vérifie $\lnot P(x_0)$, alors $x_0$ est un contre-exemple et
l'\omterm{def:b1:logic:statement}{assertion} universelle est en défaut. Pour la seconde règle~: «~aucun
$x$ ne vérifie $P$~» signifie que l'\omterm{def:b1:logic:sets}{ensemble} $\{x : P(x)\}$ est vide,
c'est-à-dire que tout $x$ appartient à son complémentaire $A$.
Appliquées en cascade à un préfixe de quantificateurs imbriqués, les
deux règles donnent le procédé mécanique de l'\cref{ex:b1:logic:limit}~:
la négation parcourt la phrase de gauche à droite, échangeant chaque
$\forall$ en $\exists$ et chaque $\exists$ en $\forall$, et nie
finalement le prédicat le plus intérieur.
\end{proof}

\begin{example}[Nier des phrases mathématiques courantes]\label{ex:b1:logic:negpractice}
Soit $f \colon \R \to \R$. La phrase «~$f$ est croissante~» s'écrit
\[
\forall x \in \R,\ \forall y \in \R,\quad
x \leq y \implies f(x) \leq f(y) ,
\]
et sa négation, d'après la \cref{prop:b1:logic:negquant} et la règle
$\lnot(P \implies Q) \iff P \land \lnot Q$~:
\[
\exists x \in \R,\ \exists y \in \R,\quad
x \leq y \ \text{ et }\ f(x) > f(y) :
\]
un seul couple témoin suffit. De même, «~$f$ est bornée~» s'écrit
$\exists M \in \R,\ \forall x \in \R,\ \abs{f(x)} \leq M$, de négation
\[
\forall M \in \R,\ \exists x \in \R,\quad \abs{f(x)} > M :
\]
\emph{quelle que soit} la borne proposée, un point la dépasse. L'idée à
retenir~: une négation correcte ne contient jamais de «~non~» appliqué à
un bloc quantifié --- c'est une nouvelle \omterm{def:b1:logic:statement}{assertion} positive, dans
laquelle les rôles sont échangés~: on produit désormais les témoins
qu'on recevait auparavant.
\end{example}

\begin{example}[Ordre des quantificateurs]\label{ex:b1:logic:order}
L'\omterm{def:b1:logic:order}{ordre} de deux quantificateurs différents compte~:
\[
\forall x \in \R,\ \exists y \in \R,\ y > x
\quad\text{est vraie (prendre } y = x+1\text{),}
\]
\[
\exists y \in \R,\ \forall x \in \R,\ y > x
\quad\text{est fausse (aucun réel ne dépasse tous les réels).}
\]
Dans la première \omterm{def:b1:logic:statement}{assertion}, $y$ peut dépendre de $x$~; dans la seconde,
un seul et même $y$ doit convenir pour tous les $x$. Deux
quantificateurs identiques, en revanche, commutent toujours.
\end{example}

\begin{example}[Lire une définition à trois quantificateurs]\label{ex:b1:logic:limit}
La phrase «~la suite $(u_n)$ converge vers $\ell$~» s'écrira dans le
\cref{ch:b1:seq} sous la forme
\[
\forall \varepsilon > 0,\ \exists N \in \N,\ \forall n \geq N,\quad
\abs{u_n - \ell} \leq \varepsilon .
\]
Sa négation, par la \cref{prop:b1:logic:negquant} appliquée trois fois, est
\[
\exists \varepsilon > 0,\ \forall N \in \N,\ \exists n \geq N,\quad
\abs{u_n - \ell} > \varepsilon .
\]
Savoir nier de telles phrases mécaniquement, sans réfléchir à ce
qu'elles signifient, est une véritable compétence~: elle sépare le
travail logique du travail mathématique.
\end{example}

\section{Techniques de démonstration}

\begin{method}[Les schémas de démonstration usuels]\label{met:b1:logic:proofs}
Pour démontrer\dots
\begin{enumerate}
  \item \emph{une implication $P \implies Q$ directement}~: supposer
        $P$, en déduire $Q$~;
  \item \emph{par contraposition}~: supposer $\lnot Q$, en déduire
        $\lnot P$ --- licite d'après la \cref{prop:b1:logic:rules}~(4)~;
  \item \emph{par l'absurde}~: supposer l'\omterm{def:b1:logic:statement}{assertion} fausse, en tirer
        une contradiction~;
  \item \emph{une équivalence}~: démontrer les deux implications
        séparément (ou enchaîner des équivalences connues)~;
  \item \emph{une \omterm{def:b1:logic:statement}{assertion} «~pour tout~»}~: se donner un $x$
        \emph{quelconque} de $E$ («~soit $x \in E$~») et démontrer
        $P(x)$~;
  \item \emph{une \omterm{def:b1:logic:statement}{assertion} «~il existe~»}~: exhiber un témoin, ou
        démontrer l'existence indirectement~;
  \item \emph{par récurrence}~: voir \cref{thm:b1:logic:induction}.
\end{enumerate}
Lorsqu'on démontre une \omterm{def:b1:logic:statement}{assertion} portant sur un élément quelconque mais
bien choisi, il ne faut jamais lui prêter de propriétés supplémentaires~:
«~soit $x \in \R$~» suivi de «~comme $x > 0$\dots~» ne démontre rien sur
les $x$ négatifs.
\end{method}

\begin{remark}[Pièges classiques des démonstrations]\label{rem:b1:logic:pitfalls}
Quatre traquenards classiques, qu'il vaut la peine de nommer une fois.
\begin{enumerate}
  \item \emph{La réciproque à la place de la contraposée.} $Q \implies
        P$ n'est \emph{pas} équivalente à $P \implies Q$~; seule
        $\lnot Q \implies \lnot P$ l'est. «~S'il pleut, la rue est
        mouillée~» n'autorise pas à conclure qu'il pleut d'une rue
        mouillée.
  \item \emph{Démontrer une équivalence par une seule implication.} Une
        affirmation «~si et seulement si~» est deux théorèmes~;
        annoncer quel sens on démontre, et démontrer les deux. Les
        chaînes de $\iff$ ne sont licites que si \emph{chaque} maillon
        est réellement réversible --- élever une équation au carré, par
        exemple, ne l'est pas.
  \item \emph{Les démonstrations à l'envers.} Partir de la conclusion
        souhaitée et en déduire une \omterm{def:b1:logic:statement}{assertion} vraie ne démontre rien
        (de $-1 = 1$ on tire le vrai $1 = 1$ en élevant au carré). Un
        calcul peut être \emph{découvert} à l'envers, mais il doit être
        \emph{écrit} à l'endroit, ou avec des équivalences explicites.
  \item \emph{Témoin fixé contre élément quelconque.} Pour démontrer
        $\exists x,\ P(x)$, on peut exhiber un seul $x$ habilement
        choisi~; pour démontrer $\forall x,\ P(x)$, le $x$ choisi doit
        rester quelconque. Confondre les deux --- vérifier une
        \omterm{def:b1:logic:statement}{assertion} universelle sur un exemple --- est l'erreur la plus
        fréquente dans les copies de débutants.
\end{enumerate}
\end{remark}

\begin{example}[Contraposition et absurde à l'œuvre]\label{ex:b1:logic:sqrt2}
\emph{Pour $n \in \N$~: si $n^2$ est pair alors $n$ est pair.} Par
contraposition~: si $n$ est impair, $n = 2k+1$, alors
$n^2 = 4k^2 + 4k + 1$ est impair.

\emph{$\sqrt 2$ est irrationnel.} Par l'absurde~: supposons
$\sqrt 2 = p/q$ avec $p, q \in \N^*$ et la fraction irréductible.
Alors $p^2 = 2q^2$ est pair, donc $p$ est pair (point précédent),
$p = 2r$~; alors $q^2 = 2r^2$ est pair, donc $q$ est pair --- ce qui
contredit l'irréductibilité.
\end{example}

\begin{theorem}[Récurrence]\label{thm:b1:logic:induction}
Soit $P(n)$ une propriété de l'entier $n$. Si
\begin{enumerate}
  \item $P(0)$ est vraie, et
  \item pour tout $n \in \N$, $P(n) \implies P(n+1)$,
\end{enumerate}
alors $P(n)$ est vraie pour tout $n \in \N$.

\emph{Récurrence forte~:} la conclusion est inchangée si (2) est
remplacée par~: pour tout $n$, $\bigl(P(0) \land \dots \land
P(n)\bigr) \implies P(n+1)$.
\end{theorem}

\begin{proof}
C'est une propriété de $\N$ lui-même, équivalente à~: \emph{toute partie
non vide de $\N$ admet un plus petit élément} (que nous admettons
connue). En effet, supposons (1) et (2) vérifiées et posons
$A = \{n \in \N : P(n) \text{ fausse}\}$. Si $A \neq \emptyset$, il
admet un plus petit élément $m$~; $m \neq 0$ d'après (1)~; alors
$m - 1 \notin A$, donc $P(m-1)$ est vraie, et (2) donne $P(m)$ ---
contradiction. Donc $A = \emptyset$. Pour la récurrence forte, le même
argument s'applique~: $P(0), \dots, P(m-1)$ sont toutes vraies puisque
$m$ est le plus petit élément de $A$.
\end{proof}

\begin{example}[Démontrer une existence et unicité]\label{ex:b1:logic:uniqueexistence}
Une \omterm{def:b1:logic:statement}{assertion} $\exists!\,x,\ P(x)$ est \emph{deux} \omterm{def:b1:logic:statement}{assertions},
démontrées séparément~: l'existence (exhiber ou construire un $x_0$
tel que $P(x_0)$) et l'unicité (supposer $P(x)$ et $P(x')$, en
déduire $x = x'$). Exemple~: \emph{il existe un unique réel $x$ tel
que $x^3 + x = 2$.} Existence~: $x_0 = 1$ convient, puisque
$1 + 1 = 2$. Unicité~: si $x^3 + x = x'^3 + x'$, alors
\[
0 = (x^3 - x'^3) + (x - x')
= (x - x')\,\bigl(x^2 + xx' + x'^2 + 1\bigr),
\]
et le second facteur est strictement positif (il vaut $\bigl(x +
\tfrac{x'}2\bigr)^2 + \tfrac34 x'^2 + 1 \geq 1$), donc $x = x'$.
Remarquons le partage des tâches~: l'existence a reposé sur une
intuition heureuse, l'unicité sur une algèbre valable pour des
solutions \emph{quelconques} --- aucun des deux arguments ne fait le
travail de l'autre, et oublier la seconde moitié est une tentation
permanente une fois qu'une solution a été trouvée.
\end{example}

\begin{example}\label{ex:b1:logic:induction}
Pour tout $n \in \N^*$~: $\;\sum_{k=1}^n k = \frac{n(n+1)}{2}$.
Initialisation $n = 1$~: les deux membres valent $1$. Hérédité~: en
supposant la formule au rang $n$,
\[
\sum_{k=1}^{n+1} k = \frac{n(n+1)}{2} + (n+1)
= (n+1)\Bigl(\frac n2 + 1\Bigr) = \frac{(n+1)(n+2)}{2}. \qedhere
\]
\end{example}

\begin{example}[La récurrence forte à l'œuvre]\label{ex:b1:logic:strongind}
\emph{Tout entier $n \geq 2$ est un produit de nombres premiers} (un
nombre premier étant un entier $\geq 2$ dont les seuls diviseurs
$\geq 1$ sont $1$ et lui-même~; les nombres premiers sont étudiés pour
eux-mêmes au \cref{ch:b1:arith}). La récurrence ordinaire est ici
impuissante~: savoir que $95 = 5 \times 19$ se factorise ne dit rien de
$96$. La récurrence forte convient exactement. Initialisation~: $2$ est
premier, donc produit (à un facteur) de nombres premiers. Hérédité~:
soit $n \geq 2$ et supposons que tout entier $m$ avec $2 \leq m \leq n$
est un produit de nombres premiers. Si $n + 1$ est premier, c'est fini.
Sinon $n + 1 = ab$ avec $2 \leq a, b \leq n$~; par l'hypothèse forte,
$a$ et $b$ sont tous deux des produits de nombres premiers, donc $n + 1$
aussi. L'idée à retenir~: la récurrence forte est l'outil adapté chaque
fois que la «~raison~» de $P(n+1)$ se trouve à un rang antérieur
imprévisible, et non au rang $n$.
\end{example}

\section{Ensembles}

\begin{definition}[Opérations sur les ensembles]\label{def:b1:logic:sets}
Nous prenons la notion d'\emph{ensemble}\index{ensemble} et la relation
d'appartenance $x \in E$ comme primitives. Pour des ensembles $A, B$
inclus dans un ensemble ambiant $E$~:
\begin{itemize}
  \item l'\emph{inclusion}~: $A \subseteq B$ lorsque $\forall x,\ x \in
        A \implies x \in B$~; l'égalité $A = B$ lorsque $A \subseteq B$
        et $B \subseteq A$~;
  \item la \emph{réunion} $A \cup B$, l'\emph{intersection} $A \cap B$,
        la \emph{différence} $A \setminus B = \{x \in A : x \notin B\}$,
        le \emph{complémentaire} $\overline{A} = E \setminus A$~;
  \item l'\emph{ensemble vide} $\emptyset$, contenu dans tout ensemble~;
  \item l'\emph{ensemble des parties}\index{ensemble des parties}
        $\mathcal{P}(E)$~: l'ensemble de toutes les parties de $E$~;
  \item le \emph{produit} $E \times F$~: l'ensemble des couples
        $(x, y)$ avec $x \in E$, $y \in F$.
\end{itemize}
\end{definition}

\begin{example}[Se familiariser avec l'ensemble des parties]\label{ex:b1:logic:powerset}
Pour $E = \{a, b\}$~:
\[
\mathcal P(E) = \bigl\{\, \emptyset,\ \{a\},\ \{b\},\ \{a, b\}
\,\bigr\},
\]
quatre éléments --- et notons la discipline des types~: $a \in E$ mais
$\{a\} \in \mathcal P(E)$~; les \omterm{def:b1:logic:statement}{assertions} $a \in \mathcal P(E)$ et
$\{a\} \subseteq \mathcal P(E)$ sont toutes deux fausses telles
qu'elles sont écrites (la seconde exigerait que $a$ soit une
\emph{partie} de $E$). En itérant à partir de rien~:
$\mathcal P(\emptyset) = \{\emptyset\}$ a un élément,
$\mathcal P(\mathcal P(\emptyset)) = \{\emptyset, \{\emptyset\}\}$ en a
deux, le suivant en a quatre --- les \omterm{def:b1:logic:sets}{ensembles} d'\omterm{def:b1:logic:sets}{ensembles} sont des
\omterm{def:b1:logic:sets}{ensembles} ordinaires, et le \cref{ch:b1:counting} confirmera le
doublement~: $\abs{\mathcal P(E)} = 2^{\abs E}$. Garder les niveaux
($x$, $\{x\}$, $\{\{x\}\}$) bien distincts, c'est déjà la moitié du
travail dans des exercices comme
les \cref{exo:b1:logic:11,exo:b1:logic:12}.
\end{example}

\begin{proposition}[Algèbre des ensembles]\label{prop:b1:logic:setalgebra}
Pour des parties $A, B, C$ de $E$~:
\begin{enumerate}
  \item $A \cap (B \cup C) = (A \cap B) \cup (A \cap C)$ et
        $A \cup (B \cap C) = (A \cup B) \cap (A \cup C)$~;
  \item De Morgan~: $\overline{A \cup B} = \overline{A} \cap
        \overline{B}$ et $\overline{A \cap B} = \overline{A} \cup
        \overline{B}$~;
  \item $A \subseteq B \iff \overline{B} \subseteq \overline{A}$.
\end{enumerate}
\end{proposition}

\begin{proof}
Chaque identité traduit une règle de la \cref{prop:b1:logic:rules} au
moyen du dictionnaire (appartenir ou non à $A$) $\leftrightarrow$
(\omterm{def:b1:logic:statement}{assertion} vraie ou fausse)~: par exemple $x \in \overline{A \cup B}
\iff \lnot(x \in A \lor x \in B) \iff (x \notin A) \land (x \notin B)
\iff x \in \overline{A} \cap \overline{B}$. Le point (3) est la
contraposition. Comme deuxième échantillon, la première loi de
distributivité en entier~:
\[
x \in A \cap (B \cup C)
\iff (x \in A) \land \bigl(x \in B \lor x \in C\bigr)
\iff \bigl(x \in A \land x \in B\bigr) \lor
\bigl(x \in A \land x \in C\bigr),
\]
par la distributivité de la \cref{prop:b1:logic:rules}~(6), et la dernière
\omterm{def:b1:logic:statement}{assertion} se lit $x \in (A \cap B) \cup (A \cap C)$. Toute identité
ensembliste de ce genre se démontre par cette unique traduction
mécanique --- raison pour laquelle aucune n'a besoin d'être apprise par
cœur.
\end{proof}

\begin{method}[Démontrer une égalité d'ensembles]\label{met:b1:logic:doubleinc}
Pour démontrer $A = B$, on démontre les deux inclusions~: soit $x \in
A$, on montre $x \in B$~; puis soit $x \in B$, on montre $x \in A$.
Autre possibilité~: enchaîner des équivalences $x \in A \iff \dots \iff
x \in B$ lorsque chaque étape en est réellement une.
\end{method}

\begin{omfigure}
\begin{tikzpicture}[scale=0.9]
  \draw[thick] (-2.6,-1.9) rectangle (2.6,1.9);
  \node[below right] at (-2.6,1.9) {$E$};
  \begin{scope}
    \clip (-2.6,-1.9) rectangle (2.6,1.9);
    \fill[omDef!20] (-2.6,-1.9) rectangle (2.6,1.9);
    \fill[white] (-0.7,0) circle (1.1);
    \fill[white] (0.7,0) circle (1.1);
  \end{scope}
  \draw[omDef, thick] (-0.7,0) circle (1.1);
  \draw[omDef, thick] (0.7,0) circle (1.1);
  \node[omDef] at (-1.35,0.5) {$A$};
  \node[omDef] at (1.35,0.5) {$B$};
  \begin{scope}[xshift=6.6cm]
    \draw[thick] (-2.6,-1.9) rectangle (2.6,1.9);
    \node[below right] at (-2.6,1.9) {$E$};
    \begin{scope}
      \clip (-2.6,-1.9) rectangle (2.6,1.9);
      \fill[omThm!20] (-2.6,-1.9) rectangle (2.6,1.9);
      \begin{scope}
        \clip (-0.7,0) circle (1.1);
        \fill[white] (0.7,0) circle (1.1);
      \end{scope}
    \end{scope}
    \draw[omThm, thick] (-0.7,0) circle (1.1);
    \draw[omThm, thick] (0.7,0) circle (1.1);
    \node[omThm] at (-1.35,0.5) {$A$};
    \node[omThm] at (1.35,0.5) {$B$};
  \end{scope}
\end{tikzpicture}

{\small Les \omterm{prop:b1:logic:rules}{lois de De Morgan} en images~: la région grisée de gauche
est $\overline{A \cup B} = \overline A \cap \overline B$ (tout ce qui
est hors des deux disques)~; à droite, $\overline{A \cap B} = \overline
A \cup \overline B$ (tout sauf la lentille commune). Un dessin n'est
pas une démonstration, mais il rend la démonstration par éléments de
la \cref{prop:b1:logic:setalgebra} impossible à mal retenir.}
\end{omfigure}

\section{Applications}

\begin{definition}[Application, image, image réciproque]\label{def:b1:logic:map}
Une \emph{application}\index{application} (ou \emph{fonction})
$f \colon E \to F$ associe à tout élément $x$ de l'\omterm{def:b1:logic:sets}{ensemble} $E$
(l'\emph{\omterm{def:b1:logic:sets}{ensemble} de départ}) exactement un élément $f(x)$ de
l'\omterm{def:b1:logic:sets}{ensemble} $F$ (l'\emph{\omterm{def:b1:logic:sets}{ensemble} d'arrivée}). Pour $A \subseteq E$ et
$B \subseteq F$~:
\[
f(A) = \{f(x) : x \in A\} \subseteq F,
\qquad
f^{-1}(B) = \{x \in E : f(x) \in B\} \subseteq E
\]
sont l'\emph{image directe} de $A$ et l'\emph{image réciproque}
\index{image réciproque} de $B$. La \emph{composée} de $f \colon E \to
F$ et $g \colon F \to G$ est $g \circ f \colon E \to G$, $x \mapsto
g(f(x))$.
\end{definition}

\begin{remark}
La notation $f^{-1}(B)$ ne présuppose \emph{pas} l'existence d'une
\omterm{def:b1:logic:map}{application} réciproque~: $f^{-1}(B)$ est définie pour toute $f$. Les
\omterm{def:b1:logic:map}{images réciproques} se comportent mieux que les images directes~:
$f^{-1}$ préserve réunions, intersections et complémentaires, alors que
$f(A \cap A') \subseteq f(A) \cap f(A')$ peut être stricte
(\cref{exo:b1:logic:8}).
\end{remark}

\begin{example}[Calculer images et images réciproques]\label{ex:b1:logic:images}
Soit $f \colon \R \to \R$, $x \mapsto x^2$. Alors~:
\[
f\bigl(\intcc{-1}{2}\bigr) = \intcc04, \qquad
f^{-1}\bigl(\intcc14\bigr) = \intcc{-2}{-1} \cup \intcc12, \qquad
f^{-1}(\{-1\}) = \emptyset .
\]
Pour la première~: tout $x \in \intcc{-1}2$ vérifie $x^2 \in \intcc04$,
et tout $y \in \intcc04$ est atteint sous la forme $y = (\sqrt y)^2$
avec $\sqrt y \in \intcc02 \subseteq \intcc{-1}2$ --- notons que l'image
n'est \emph{pas} $\intcc14 = \{(-1)^2, 2^2\}$~: les images d'intervalles
ne se calculent pas à partir des seules extrémités. Pour la deuxième~:
$1 \leq x^2 \leq 4 \iff 1 \leq \abs x \leq 2$, ce qui se scinde en deux
morceaux. La troisième illustre qu'une \omterm{def:b1:logic:map}{image réciproque} peut être vide
--- $f^{-1}(B)$ a toujours un sens, si petite que soit l'intersection de
$B$ avec l'image. Observons enfin sur cet exemple le phénomène de
stricte inclusion de la remarque précédente~: avec $A = \intcc{-1}0$ et
$A' = \intcc01$, on a $f(A \cap A') = f(\{0\}) = \{0\}$, tandis que
$f(A) \cap f(A') = \intcc01$.
\end{example}

\begin{definition}[Injective, surjective, bijective]\label{def:b1:logic:inj}
Une \omterm{def:b1:logic:map}{application} $f \colon E \to F$ est~:
\begin{itemize}
  \item \emph{injective}\index{injective} lorsque des éléments distincts
        ont des images distinctes~: $\forall x, x' \in E,\ f(x) = f(x')
        \implies x = x'$~;
  \item \emph{surjective}\index{surjective} lorsque tout élément de $F$
        est atteint~: $\forall y \in F,\ \exists x \in E,\ f(x) = y$~;
  \item \emph{bijective}\index{bijective} lorsqu'elle est les deux,
        c'est-à-dire lorsque tout $y \in F$ a exactement un antécédent.
\end{itemize}
\end{definition}

\begin{theorem}[Application réciproque]\label{thm:b1:logic:inverse}
Une \omterm{def:b1:logic:map}{application} $f \colon E \to F$ est \omterm{def:b1:logic:inj}{bijective} si et seulement s'il
existe une \omterm{def:b1:logic:map}{application} $g \colon F \to E$ telle que $g \circ f =
\mathrm{id}_E$ et $f \circ g = \mathrm{id}_F$. Dans ce cas $g$ est
unique~; on la note $f^{-1}$ et on l'appelle la \emph{réciproque} de
$f$, et $f^{-1}$ est elle-même \omterm{def:b1:logic:inj}{bijective}, avec $(f^{-1})^{-1} = f$.
\end{theorem}

\begin{proof}
($\Rightarrow$) Si $f$ est \omterm{def:b1:logic:inj}{bijective}, tout $y \in F$ a un unique
antécédent~; définissons $g(y)$ comme cet antécédent. Alors
$f(g(y)) = y$ par construction, et $g(f(x)) = x$ parce que $x$ est
\emph{l'}antécédent de $f(x)$.

($\Leftarrow$) Supposons qu'une telle $g$ existe. Si $f(x) = f(x')$, en
appliquant $g$ on obtient $x = x'$~: $f$ est \omterm{def:b1:logic:inj}{injective}. Pour $y \in F$,
$x = g(y)$ vérifie $f(x) = y$~: $f$ est \omterm{def:b1:logic:inj}{surjective}.

Unicité~: si $g$ et $h$ conviennent toutes deux, alors $g = g \circ
\mathrm{id}_F = g \circ (f \circ h) = (g \circ f) \circ h = h$. Enfin
le couple d'identités est symétrique en $f$ et $g$, donc $g = f^{-1}$
est \omterm{def:b1:logic:inj}{bijective}, de réciproque $f$.
\end{proof}

\begin{example}[Calculer une réciproque en pratique]\label{ex:b1:logic:inversecomp}
Soit $f \colon \R \to \intoo0{+\infty}$, $f(x) = \eu^{2x+1}$. Pour
l'inverser, résolvons $y = f(x)$ pour un $y > 0$ donné~:
\[
y = \eu^{2x+1} \iff \ln y = 2x + 1 \iff x = \frac{\ln y - 1}2 ,
\]
chaque étape étant réversible sur les \omterm{def:b1:logic:sets}{ensembles} annoncés. Le calcul
livre tout d'un coup~: pour chaque $y$ de l'\omterm{def:b1:logic:sets}{ensemble} d'arrivée il y a
exactement une solution $x$, donc $f$ est \omterm{def:b1:logic:inj}{bijective}, et
\[
f^{-1} \colon \intoo0{+\infty} \to \R,
\qquad
f^{-1}(y) = \frac{\ln y - 1}2 .
\]
Une vérification rapide des deux composées ($f^{-1}(f(x)) =
\frac{(2x+1) - 1}2 = x$ et $f(f^{-1}(y)) = \eu^{\ln y} = y$) confirme
le critère du \cref{thm:b1:logic:inverse}. L'idée à retenir~:
«~résoudre en $x$ en surveillant les équivalences~» est à la fois la
démonstration d'existence, celle d'unicité et la formule --- mais cela
ne fonctionne que si l'\omterm{def:b1:logic:sets}{ensemble} d'arrivée a été correctement annoncé
($f$ n'est \emph{pas} \omterm{def:b1:logic:inj}{surjective} sur $\R$).
\end{example}

\begin{proposition}[Composition et les trois propriétés]\label{prop:b1:logic:comp}
Soient $f \colon E \to F$ et $g \colon F \to G$.
\begin{enumerate}
  \item Si $f$ et $g$ sont \omterm{def:b1:logic:inj}{injectives} (resp.\ \omterm{def:b1:logic:inj}{surjectives}, \omterm{def:b1:logic:inj}{bijectives}),
        alors $g \circ f$ l'est aussi~; et dans le cas \omterm{def:b1:logic:inj}{bijectif},
        $(g \circ f)^{-1} = f^{-1} \circ g^{-1}$.
  \item Si $g \circ f$ est \omterm{def:b1:logic:inj}{injective}, alors $f$ est \omterm{def:b1:logic:inj}{injective}. Si
        $g \circ f$ est \omterm{def:b1:logic:inj}{surjective}, alors $g$ est \omterm{def:b1:logic:inj}{surjective}.
\end{enumerate}
\end{proposition}

\begin{proof}
(1) Si $g(f(x)) = g(f(x'))$, l'\omterm{def:b1:logic:inj}{injectivité} de $g$ donne $f(x) = f(x')$,
puis l'\omterm{def:b1:logic:inj}{injectivité} de $f$ donne $x = x'$. Si $z \in G$, la \omterm{def:b1:logic:inj}{surjectivité}
de $g$ fournit $y$ avec $g(y) = z$, puis la \omterm{def:b1:logic:inj}{surjectivité} de $f$ fournit
$x$ avec $f(x) = y$, de sorte que $g(f(x)) = z$. Dans le cas \omterm{def:b1:logic:inj}{bijectif},
on vérifie directement que $f^{-1} \circ g^{-1}$ est un inverse à
gauche et à droite de $g \circ f$, et l'unicité du
\cref{thm:b1:logic:inverse} conclut.

(2) Si $f(x) = f(x')$ alors $g(f(x)) = g(f(x'))$, et l'\omterm{def:b1:logic:inj}{injectivité} de
$g \circ f$ donne $x = x'$. Si $z \in G$, la \omterm{def:b1:logic:inj}{surjectivité} de $g \circ f$
fournit $x$ avec $g(f(x)) = z$~: alors $y = f(x)$ vérifie $g(y) = z$.
\end{proof}

\begin{example}[Le point (2) est optimal]\label{ex:b1:logic:compsharp}
Dans la \cref{prop:b1:logic:comp}~(2), on ne peut pas renforcer les
conclusions~: $g \circ f$ \omterm{def:b1:logic:inj}{bijective} n'oblige \emph{pas} $f$ à être
\omterm{def:b1:logic:inj}{surjective} ni $g$ à être \omterm{def:b1:logic:inj}{injective}. Prenons $E = G = \{1\}$,
$F = \{1, 2\}$, avec $f(1) = 1$ et $g(1) = g(2) = 1$~: alors
$g \circ f = \mathrm{id}_E$ est \omterm{def:b1:logic:inj}{bijective}, et pourtant $f$ manque
l'élément $2$ et $g$ écrase les deux éléments. La morale est une règle
de comptabilité précise~: l'information de la composée descend vers
l'\omterm{def:b1:logic:map}{application} \emph{intérieure} pour l'\omterm{def:b1:logic:inj}{injectivité} et vers
l'\omterm{def:b1:logic:map}{application} \emph{extérieure} pour la \omterm{def:b1:logic:inj}{surjectivité}, jamais dans
l'autre sens. (L'\cref{exo:b1:logic:9} construit le même phénomène avec
des \omterm{def:b1:logic:sets}{ensembles} infinis, où il est le moteur des inverses d'un seul
côté.)
\end{example}

\begin{example}\label{ex:b1:logic:injexamples}
$f \colon \R \to \R$, $x \mapsto x^2$ n'est ni \omterm{def:b1:logic:inj}{injective} ($f(-1) =
f(1)$) ni \omterm{def:b1:logic:inj}{surjective} ($-1$ n'a pas d'antécédent). En restreignant
l'\omterm{def:b1:logic:sets}{ensemble} de départ et l'\omterm{def:b1:logic:sets}{ensemble} d'arrivée, $f \colon \R_+ \to \R_+$,
$x \mapsto x^2$ est \omterm{def:b1:logic:inj}{bijective}, de réciproque $y \mapsto \sqrt y$.
L'\omterm{def:b1:logic:inj}{injectivité} ou la \omterm{def:b1:logic:inj}{surjectivité} d'une \omterm{def:b1:logic:map}{application} dépend des \omterm{def:b1:logic:sets}{ensembles}
de départ et d'arrivée annoncés, et pas seulement de la formule.
\end{example}

\section{Relations}

\begin{definition}[Relation d'équivalence]\label{def:b1:logic:equiv}
Une \emph{relation binaire} $\mathcal{R}$ sur un \omterm{def:b1:logic:sets}{ensemble} $E$ est une
\emph{relation d'équivalence}\index{relation d'équivalence} lorsqu'elle
est~: \emph{réflexive} ($x \mathbin{\mathcal{R}} x$ pour tout $x$),
\emph{symétrique} ($x \mathbin{\mathcal{R}} y \implies y
\mathbin{\mathcal{R}} x$) et \emph{transitive} ($x
\mathbin{\mathcal{R}} y$ et $y \mathbin{\mathcal{R}} z$ entraînent $x
\mathbin{\mathcal{R}} z$). La \emph{classe d'équivalence} de $x$ est
$\mathrm{cl}(x) = \{y \in E : x \mathbin{\mathcal{R}} y\}$.
\end{definition}

\begin{example}[Vérifier les trois axiomes]\label{ex:b1:logic:fracpart}
Sur $\R$, déclarons $x \mathbin{\mathcal{R}} y$ lorsque $x - y \in \Z$.
\emph{Réflexive~:} $x - x = 0 \in \Z$. \emph{Symétrique~:} si $x - y
\in \Z$ alors $y - x = -(x - y) \in \Z$. \emph{Transitive~:} si $x -
y \in \Z$ et $y - z \in \Z$, alors $x - z = (x - y) + (y - z) \in
\Z$ (somme d'entiers). Donc $\mathcal R$ est une \omterm{def:b1:logic:equiv}{relation
d'équivalence}, et $\mathrm{cl}(x) = x + \Z = \{x + k : k \in \Z\}$~:
chaque classe contient exactement un représentant dans $\intco01$, sa
\emph{partie fractionnaire}. En revanche, la relation «~$\abs{x - y}
\leq 1$~» sur $\R$ est réflexive et symétrique mais \emph{pas}
transitive ($0 \mathbin{\mathcal R} 1$ et $1 \mathbin{\mathcal R}
2$, alors que $\abs{0 - 2} > 1$)~: la proximité ne se propage pas, et
aucune \omterm{thm:b1:logic:partition}{partition} en classes n'existe --- un contre-exemple utile à
garder en tête quand la vérification des axiomes semble routinière.
\end{example}

\begin{theorem}[Les classes forment une partition]\label{thm:b1:logic:partition}
Soit $\mathcal{R}$ une \omterm{def:b1:logic:equiv}{relation d'équivalence} sur $E$. Alors les classes
d'équivalence sont non vides, deux à deux disjointes ou égales, et leur
réunion est $E$~: elles forment une \emph{partition}\index{partition} de
$E$. Réciproquement, toute \omterm{thm:b1:logic:partition}{partition} de $E$ provient ainsi d'exactement
une \omterm{def:b1:logic:equiv}{relation d'équivalence} («~être dans le même morceau~»).
\end{theorem}

\begin{proof}
$x \in \mathrm{cl}(x)$ par réflexivité, donc les classes sont non vides
et de réunion $E$. Supposons $\mathrm{cl}(x) \cap \mathrm{cl}(y) \neq
\emptyset$, disons que $z$ appartient aux deux. Alors $x
\mathbin{\mathcal{R}} z$ et $y \mathbin{\mathcal{R}} z$, donc par
symétrie et transitivité $x \mathbin{\mathcal{R}} y$. Alors, pour tout
$t \in \mathrm{cl}(y)$, la transitivité donne $t \in \mathrm{cl}(x)$, et
symétriquement~: les deux classes sont égales. Pour la réciproque, soit
$(E_i)_{i \in I}$ une \omterm{thm:b1:logic:partition}{partition} de $E$ et définissons
$x \mathbin{\mathcal S} y$ par «~un morceau contient à la fois $x$ et
$y$~». \emph{Réflexive~:} $x$ appartient à un morceau, qui le contient
alors deux fois. \emph{Symétrique~:} la condition définissante est
symétrique en $x$ et $y$. \emph{Transitive~:} si $x, y \in E_i$ et
$y, z \in E_j$, alors $y \in E_i \cap E_j$, donc $E_i = E_j$ (des
morceaux distincts sont disjoints) et $x, z$ partagent un morceau. La
classe de $x$ pour $\mathcal S$ est exactement le morceau contenant
$x$, donc les classes sont les morceaux donnés. Enfin la relation est
déterminée par ses classes~: deux \omterm{def:b1:logic:equiv}{relations d'équivalence} ayant les
mêmes classes relient les mêmes couples, puisque chacune relie $x$ et
$y$ exactement quand $y$ appartient à la classe de $x$ --- d'où
l'unicité annoncée.
\end{proof}

\begin{example}\label{ex:b1:logic:congruence}
Sur $\Z$, la congruence modulo $n$ ($x \equiv y \pmod n$ lorsque $n$
divise $x - y$) est une \omterm{def:b1:logic:equiv}{relation d'équivalence}~; ses classes sont les
$n$ \omterm{def:b1:logic:sets}{ensembles} d'entiers ayant un reste donné dans la division par $n$.
Cet exemple devient l'anneau $\Z/n\Z$ au \cref{ch:b1:structures}.
\end{example}

\begin{definition}[Relation d'ordre]\label{def:b1:logic:order}
Une relation $\preceq$ sur $E$ est un
\emph{ordre}\index{relation d'ordre} lorsqu'elle est réflexive,
\emph{antisymétrique} ($x \preceq y$ et $y \preceq x$ entraînent
$x = y$) et transitive. L'ordre est
\emph{total} lorsque deux éléments quelconques sont comparables,
\emph{partiel} sinon. Un élément $M \in A \subseteq E$ est un
\emph{plus grand élément} de $A$ lorsque $a \preceq M$ pour tout
$a \in A$~; les plus grands (et les plus petits) éléments sont uniques
lorsqu'ils existent.
\end{definition}

\begin{example}\label{ex:b1:logic:orders}
$(\R, \leq)$ est totalement ordonné. $(\mathcal{P}(E), \subseteq)$ est
partiellement ordonné dès que $E$ a deux éléments~: $\{a\}$ et
$\{b\}$ ne sont pas comparables. La partie $A = \{\{a\}, \{b\}\}$ de
$\mathcal{P}(\{a,b\})$ n'a pas de plus grand élément, et pourtant elle
admet le majorant $\{a, b\}$~: la distinction entre plus grand élément
et majorant reviendra, pour $\R$, au \cref{ch:b1:reals}.
\end{example}

\begin{example}[Deux ordres sur la grille $\N^2$]\label{ex:b1:logic:lexorder}
Sur les couples d'entiers naturels, comparons composante par
composante~: $(a, b) \preceq (a', b')$ lorsque $a \leq a'$ \emph{et}
$b \leq b'$ (l'\emph{\omterm{def:b1:logic:order}{ordre} produit}). C'est bien un \omterm{def:b1:logic:order}{ordre} --- chaque
axiome est hérité coordonnée par coordonnée --- mais un \omterm{def:b1:logic:order}{ordre} partiel~:
$(1, 3)$ et $(2, 0)$ sont incomparables. Comparons maintenant comme un
dictionnaire~: $(a, b) \preceq_{\mathrm{lex}} (a', b')$ lorsque
$a < a'$, ou bien $a = a'$ et $b \leq b'$ (l'\emph{\omterm{def:b1:logic:order}{ordre}
lexicographique}). La transitivité demande une vérification en deux cas
mais elle est vraie, et deux couples quelconques sont désormais
comparables~: l'\omterm{def:b1:logic:order}{ordre} est total. Les deux \omterm{def:b1:logic:order}{ordres} classent le même
\omterm{def:b1:logic:sets}{ensemble} différemment --- $(0, 100) \preceq_{\mathrm{lex}} (1, 0)$ alors
que l'\omterm{def:b1:logic:order}{ordre} produit ne dit rien --- ce qui rappelle qu'un \omterm{def:b1:logic:order}{ordre} est une
structure que l'on \emph{choisit}, et non une propriété de l'\omterm{def:b1:logic:sets}{ensemble}.
La comparaison lexicographique est aussi l'astuce standard pour
transformer plusieurs critères de tri en un seul.
\end{example}

\begin{remark}[Interlude~: la taille comme bijection]\label{rem:b1:logic:interlude}
Un thème discret de ce chapitre mérite d'être mis en lumière~: les
\omterm{def:b1:logic:inj}{bijections} sont la notion mathématique de «~même taille~». Pour les
\omterm{def:b1:logic:sets}{ensembles} finis, cela devient le calcul de dénombrement du
\cref{ch:b1:counting}, où chaque formule est secrètement une
\omterm{def:b1:logic:inj}{bijection}~; pour les \omterm{def:b1:logic:sets}{ensembles} infinis, cela devient le devoir maison
ci-dessous, où $\N$, $\Q$ et $\R$ se révèlent avoir des tailles
véritablement différentes. Le même dictionnaire réapparaît deux fois
encore dans ce volume, sous des formes raffinées~: les suites
(\cref{ch:b1:seq}) ne sont rien d'autre que des \omterm{def:b1:logic:map}{applications} $\N \to
\R$, si bien que les \omterm{def:b1:logic:statement}{assertions} sur les suites sont des \omterm{def:b1:logic:statement}{assertions} sur
un \omterm{def:b1:logic:sets}{ensemble} d'\omterm{def:b1:logic:map}{applications}~; et l'algèbre linéaire mesurera les espaces
vectoriels non pas par des \omterm{def:b1:logic:inj}{bijections} mais par des \omterm{def:b1:logic:inj}{bijections}
\emph{linéaires}, dont l'existence est gouvernée par un unique nombre,
la dimension (\cref{ch:b1:findim}). Chaque fois qu'une nouvelle notion
de «~mêmeté~» apparaît --- équipotence, isomorphisme de groupes
(\cref{ch:b1:structures}), isomorphisme linéaire --- le schéma du
\cref{thm:b1:logic:inverse} se répète~: être «~le même~», c'est être
relié par une \omterm{def:b1:logic:map}{application} inversible qui respecte la structure.
\end{remark}

\begin{remark}[Où ce chapitre sert]\label{rem:b1:logic:whereused}
Partout --- mais quelques endroits méritent d'être signalés. La
gymnastique à trois quantificateurs de l'\cref{ex:b1:logic:limit} est le
pain quotidien des \cref{ch:b1:seq,ch:b1:continuity}~: toute
démonstration de limite est une partie jouée contre un $\varepsilon$
quelconque. Les classes d'équivalence réapparaissent comme les classes
de congruence de $\Z/n\Z$ au \cref{ch:b1:structures}, où la \omterm{thm:b1:logic:partition}{partition} du
\cref{thm:b1:logic:partition} acquiert une structure algébrique propre.
Les \omterm{def:b1:logic:order}{relations d'ordre}, les majorants et les bornes supérieures
deviennent le cœur axiomatique de $\R$ au \cref{ch:b1:reals}. Les
\omterm{def:b1:logic:inj}{injections}, \omterm{def:b1:logic:inj}{surjections} et \omterm{def:b1:logic:inj}{bijections} reviennent comme les \omterm{def:b1:logic:map}{applications}
linéaires du \cref{ch:b1:linmaps}, où l'\omterm{def:b1:logic:inj}{injectivité} se teste sur un seul
vecteur (le noyau)~; et le devoir maison ci-dessous transforme la simple
notion de \omterm{def:b1:logic:inj}{bijection} en une théorie des \emph{tailles des \omterm{def:b1:logic:sets}{ensembles}
infinis}, dont les conclusions (dénombrabilité de $\Q$, non-dénombrabilité
de $\R$) refont surface aux \cref{ch:b1:reals,ch:b1:topology}.
\end{remark}

\section{Exercices}

\begin{exercise}[$\star$]\label{exo:b1:logic:1}
Écrire la négation de chaque \omterm{def:b1:logic:statement}{assertion}, sans employer le mot «~non~»~:
\begin{enumerate}
  \item $\forall x \in \R,\ \exists y \in \R,\ x + y > 0$~;
  \item $\exists x \in \R,\ \forall y \in \R,\ xy = 0$~;
  \item $\forall \varepsilon > 0,\ \exists \delta > 0,\ \forall x \in
        \R,\ \abs{x} \leq \delta \implies \abs{f(x)} \leq \varepsilon$
        (pour une \omterm{def:b1:logic:map}{application} fixée $f \colon \R \to \R$).
\end{enumerate}
Décider ensuite si les \omterm{def:b1:logic:statement}{assertions} (1) et (2) sont vraies.
\end{exercise}

\begin{exercise}[$\star$]\label{exo:b1:logic:2}
Soient $P, Q$ des \omterm{def:b1:logic:statement}{assertions}. À l'aide de tables de vérité, démontrer
que $\lnot(P \implies Q) \iff P \land (\lnot Q)$, et en déduire la
négation de~: «~si une fonction est dérivable alors elle est
continue~».
\end{exercise}

\begin{exercise}[$\star$]\label{exo:b1:logic:3}
Démontrer par contraposition~: pour $x \in \R$, si $x^3 + x \geq 2$
alors $x \geq 1$. Démontrer ensuite par l'absurde qu'il n'existe pas de
plus petit réel strictement positif.
\end{exercise}

\begin{exercise}[$\star$]\label{exo:b1:logic:4}
Démontrer par récurrence que, pour tout $n \in \N$~:
\begin{enumerate}
  \item $\sum_{k=0}^{n} 2^k = 2^{n+1} - 1$~;
  \item $4^n + 5$ est divisible par $3$.
\end{enumerate}
\end{exercise}

\begin{exercise}[$\star$]\label{exo:b1:logic:5}
Trouver la faille dans la «~démonstration~» suivante que tous les
crayons ont la même couleur. \emph{Soit $P(n)$~: «~dans tout \omterm{def:b1:logic:sets}{ensemble}
de $n$ crayons, tous les crayons ont la même couleur~». $P(1)$ est
clair. Supposons $P(n)$ et prenons $n+1$ crayons~; en retirant le
dernier, les $n$ premiers ont même couleur~; en retirant le premier,
les $n$ derniers ont même couleur~; donc les $n+1$ ont même couleur.}
\end{exercise}

\begin{exercise}[$\star$]\label{exo:b1:logic:6}
Soient $A, B, C$ des parties de $E$. Démontrer~:
\begin{enumerate}
  \item $A \setminus B = A \cap \overline{B}$~;
  \item $(A \cup B) \setminus C = (A \setminus C) \cup (B \setminus
        C)$~;
  \item $A \subseteq B \iff A \cup B = B \iff A \cap B = A$.
\end{enumerate}
\end{exercise}

\begin{exercise}[$\star\star$]\label{exo:b1:logic:7}
Pour chaque \omterm{def:b1:logic:map}{application}, décider (avec démonstration) si elle est
\omterm{def:b1:logic:inj}{injective}, \omterm{def:b1:logic:inj}{surjective}, \omterm{def:b1:logic:inj}{bijective}~:
\begin{enumerate}
  \item $f \colon \N \to \N$, $n \mapsto n + 1$~;
  \item $g \colon \Z \to \Z$, $n \mapsto n + 1$~;
  \item $h \colon \R \setminus \{1\} \to \R$, $x \mapsto
        \frac{x+1}{x-1}$.
\end{enumerate}
Pour $h$, ajuster l'\omterm{def:b1:logic:sets}{ensemble} d'arrivée pour la rendre \omterm{def:b1:logic:inj}{bijective} et
calculer la réciproque.
\end{exercise}

\begin{exercise}[$\star\star$]\label{exo:b1:logic:8}
Soient $f \colon E \to F$, $A, A' \subseteq E$ et $B, B' \subseteq
F$.
\begin{enumerate}
  \item Démontrer $f^{-1}(B \cap B') = f^{-1}(B) \cap f^{-1}(B')$ et
        $f(A \cup A') = f(A) \cup f(A')$.
  \item Démontrer $f(A \cap A') \subseteq f(A) \cap f(A')$ et donner un
        exemple où l'inclusion est stricte.
  \item Démontrer~: $f$ est \omterm{def:b1:logic:inj}{injective} si et seulement si $f(A \cap A')
        = f(A) \cap f(A')$ pour toutes parties $A, A'$.
\end{enumerate}
\end{exercise}

\begin{exercise}[$\star\star$]\label{exo:b1:logic:9}
Soient $f \colon E \to F$ et $g \colon F \to E$ telles que $g \circ f =
\mathrm{id}_E$. Démontrer que $f$ est \omterm{def:b1:logic:inj}{injective} et que $g$ est
\omterm{def:b1:logic:inj}{surjective}. Donner un exemple où ni $f$ ni $g$ n'est \omterm{def:b1:logic:inj}{bijective}.
\end{exercise}

\begin{exercise}[$\star\star$]\label{exo:b1:logic:10}
Sur $\R$, on définit $x \mathbin{\mathcal{R}} y \iff x^2 - y^2 = x - y$.
Démontrer que $\mathcal{R}$ est une \omterm{def:b1:logic:equiv}{relation d'équivalence} et décrire la
classe d'équivalence de chaque réel $x$. Quelles classes ont exactement
un élément~?
\end{exercise}

\begin{exercise}[$\star\star\star$]\label{exo:b1:logic:11}
(Cantor) Soit $E$ un \omterm{def:b1:logic:sets}{ensemble}. Démontrer qu'il n'existe pas de
\omterm{def:b1:logic:inj}{surjection} de $E$ sur $\mathcal{P}(E)$. \emph{Indication~: étant donnée
$f \colon E \to \mathcal{P}(E)$, considérer $D = \{x \in E : x \notin
f(x)\}$.}
\end{exercise}

\begin{exercise}[$\star\star\star$]\label{exo:b1:logic:12}
Soit $f \colon E \to F$ une \omterm{def:b1:logic:map}{application}. On définit $\Phi \colon
\mathcal{P}(F) \to \mathcal{P}(E)$ par $\Phi(B) = f^{-1}(B)$.
\begin{enumerate}
  \item Démontrer que $f$ est \omterm{def:b1:logic:inj}{surjective} si et seulement si $\Phi$ est
        \omterm{def:b1:logic:inj}{injective}.
  \item Démontrer que $f$ est \omterm{def:b1:logic:inj}{injective} si et seulement si $\Phi$ est
        \omterm{def:b1:logic:inj}{surjective}.
\end{enumerate}
\end{exercise}

\section{Problème~: comparer les infinis}

\begin{problem}\label{pb:b1:logic:1}
Quand deux \omterm{def:b1:logic:sets}{ensembles} ont-ils «~le même nombre d'éléments~»~? La réponse
de Cantor --- lorsqu'il existe une \omterm{def:b1:logic:inj}{bijection} entre eux --- se révèle
utilisable même pour des \omterm{def:b1:logic:sets}{ensembles} infinis, et elle scinde l'infini en
tailles véritablement différentes. Ce problème construit toute la boîte
à outils à partir des seules définitions de ce chapitre~: le théorème de
Cantor--Schröder--Bernstein (deux \omterm{def:b1:logic:inj}{injections} fabriquent une \omterm{def:b1:logic:inj}{bijection}),
la dénombrabilité de $\Q$, la non-dénombrabilité de $\R$ par l'argument
diagonal, et l'étonnante conclusion de Cantor en 1874~: \emph{les
\omterm{pb:b1:logic:1}{nombres transcendants} existent, et en masse}, sans qu'on en exhibe un
seul.
\index{ensemble dénombrable}\index{théorème de Cantor--Schroder--Bernstein}
Dans tout le problème, pour des \omterm{def:b1:logic:sets}{ensembles} $E$ et $F$, on écrit $E
\preceq F$ lorsqu'il existe une \omterm{def:b1:logic:inj}{injection} de $E$ dans $F$, et $E
\approx F$ («~$E$ et $F$ sont \emph{équipotents}~») lorsqu'il existe une
\omterm{def:b1:logic:inj}{bijection} de $E$ sur $F$.

\medskip
\noindent\textbf{Partie I --- Le vocabulaire de la comparaison.}
\begin{enumerate}
  \item Montrer que $\approx$ se comporte comme une \omterm{def:b1:logic:equiv}{relation
        d'équivalence}~: $E \approx E$~; si $E \approx F$ alors
        $F \approx E$~; si $E \approx F$ et $F \approx G$ alors
        $E \approx G$. (Citer précisément \cref{thm:b1:logic:inverse} et
        \cref{prop:b1:logic:comp}.)
  \item Montrer que $\preceq$ est transitive, et qu'une \omterm{def:b1:logic:inj}{injection}
        $f \colon E \to F$ induit toujours $E \approx f(E)$.
  \item Soit $E \neq \emptyset$. Montrer que $E \preceq F$ si et
        seulement s'il existe une \omterm{def:b1:logic:inj}{surjection} de $F$ sur $E$.
  \item Vérifier que $n \mapsto n + 1$ est une \omterm{def:b1:logic:inj}{bijection} de $\N$ sur
        $\N^* = \N \setminus \{0\}$, et que
        \[
        \sigma(n) = \frac n2 \ \ (n \text{ pair}), \qquad
        \sigma(n) = -\frac{n+1}2 \ \ (n \text{ impair})
        \]
        est une \omterm{def:b1:logic:inj}{bijection} de $\N$ sur $\Z$. Ainsi, retirer un point,
        ou doubler en passant aux négatifs, ne change pas la taille de
        $\N$.
\end{enumerate}

\medskip
\noindent\textbf{Partie II --- Le théorème de
Cantor--Schröder--Bernstein.} Soient $f \colon E \to F$ et $g \colon F
\to E$ deux \omterm{def:b1:logic:inj}{injections}. On pose
\[
C_0 = E \setminus g(F), \qquad C_{n+1} = g\bigl(f(C_n)\bigr)
\ \ (n \in \N), \qquad C = \bigcup_{n \in \N} C_n,
\]
et on définit $h \colon E \to F$ en envoyant $x \in C$ sur $f(x)$, et
$x \notin C$ sur l'unique $y \in F$ tel que $g(y) = x$.
\begin{enumerate}[resume]
  \item Vérifier que $h$ est bien définie~: si $x \notin C$ alors $x \in
        g(F)$, et l'élément $y$ tel que $g(y) = x$ est unique.
  \item Montrer que $g\bigl(f(C)\bigr) = \bigcup_{n \geq 1} C_n
        \subseteq C$. (Les images directes commutent aux réunions~:
        \cref{exo:b1:logic:8}.)
  \item Montrer que $h$ est \omterm{def:b1:logic:inj}{injective}. (Trois cas~; dans le cas mixte
        $x \in C$, $x' \notin C$, montrer que $h(x) = h(x')$
        forcerait $x' \in g(f(C)) \subseteq C$.)
  \item Montrer que $h$ est \omterm{def:b1:logic:inj}{surjective}~: étant donné $y \in F$,
        distinguer les cas $g(y) \notin C$ et $g(y) \in C_n$ pour un
        certain $n \geq 1$ (pourquoi $g(y) \in C_0$ est-il
        impossible~?), et exhiber dans chaque cas un antécédent de $y$.
  \item Conclure par le \emph{théorème de
        Cantor--Schröder--Bernstein}~: si $E \preceq F$ et $F \preceq
        E$, alors $E \approx F$. Commenter en une phrase ce qui rend
        cet énoncé non trivial.
  \item Deux \omterm{def:b1:logic:map}{applications}. (a) Montrer que $\intcc01 \approx \intoo01$.
        (b) Montrer que $\varphi(p, q) = 2^p(2q + 1) - 1$ définit une
        \omterm{def:b1:logic:inj}{bijection} de $\N \times \N$ sur $\N$ --- \omterm{def:b1:logic:inj}{injectivité} par un
        argument de parité, \omterm{def:b1:logic:inj}{surjectivité} par récurrence forte
        (\cref{thm:b1:logic:induction}). Ainsi $\N \times \N \approx
        \N$~: le plan des points entiers n'est pas plus gros que la
        droite.
\end{enumerate}

\medskip
\noindent\textbf{Partie III --- \omterm{pb:b1:logic:1}{Ensembles dénombrables}.} On dit qu'un
\omterm{def:b1:logic:sets}{ensemble} $E$ est \emph{au plus dénombrable} lorsque $E \preceq \N$, et
\emph{dénombrable} lorsque $E \approx \N$.
\begin{enumerate}[resume]
  \item Montrer que toute partie infinie $A \subseteq \N$ est
        dénombrable. (Définir $\varphi(n)$ par récurrence comme le plus
        petit élément de $A \setminus \{\varphi(0), \dots,
        \varphi(n-1)\}$~; montrer que $\varphi$ est strictement
        croissante, vérifie $\varphi(n) \geq n$, et atteint tout élément
        de $A$.)
  \item En déduire qu'un \omterm{def:b1:logic:sets}{ensemble} est au plus dénombrable si et
        seulement s'il est fini ou dénombrable, et observer que la
        question 9 fournit le raccourci~: si $E \preceq \N$ et $\N
        \preceq E$, alors $E$ est dénombrable.
  \item Montrer que si $E$ et $F$ sont au plus dénombrables, alors
        $E \times F$ l'est aussi. En déduire que $\Z \times \N^*$ est
        dénombrable.
  \item Montrer que $\Q$ est dénombrable. (Injecter $\Q$ dans $\Z
        \times \N^*$ en écrivant chaque rationnel sous forme
        irréductible avec dénominateur positif --- l'unicité de cette
        représentation est démontrée au \cref{ch:b1:arith}~; puis
        appliquer la question 12.)
  \item Montrer qu'une réunion dénombrable d'\omterm{def:b1:logic:sets}{ensembles} au plus
        dénombrables est au plus dénombrable~: si chaque $E_n$ ($n \in
        \N$) est au plus dénombrable, alors $\bigcup_{n \in \N} E_n$
        l'est aussi. (Envoyer $x$ sur le couple $(n, f_n(x))$ où $n$ est
        le \emph{plus petit} indice tel que $x \in E_n$.)
  \item Montrer que l'\omterm{def:b1:logic:sets}{ensemble des parties} \emph{finies} de $\N$ est
        dénombrable. (Envoyer une partie finie $F$ sur $\sum_{i \in F}
        2^i$~; démontrer l'\omterm{def:b1:logic:inj}{injectivité} en comparant le plus grand
        élément où deux parties finies diffèrent, à l'aide de
        $\sum_{k=0}^{m-1} 2^k = 2^m - 1$ tiré de
        l'\cref{exo:b1:logic:4}.)
\end{enumerate}

\medskip
\noindent\textbf{Partie IV --- Diagonalisation.} On note $\{0,1\}^{\N}$
l'\omterm{def:b1:logic:sets}{ensemble} de toutes les \omterm{def:b1:logic:map}{applications} $u \colon \N \to \{0, 1\}$,
c'est-à-dire l'\omterm{def:b1:logic:sets}{ensemble} des suites binaires.
\begin{enumerate}[resume]
  \item Construire une \omterm{def:b1:logic:inj}{bijection} entre $\mathcal{P}(\N)$ et
        $\{0,1\}^{\N}$ (fonctions indicatrices).
  \item (L'argument diagonal) Soit $\Phi \colon \N \to
        \{0,1\}^{\N}$ une \omterm{def:b1:logic:map}{application} quelconque. Considérer la suite
        $d$ définie par $d(n) = 1 - \Phi(n)(n)$. Montrer que $d$ n'est
        pas dans l'image de $\Phi$, et en conclure que $\{0,1\}^{\N}$
        n'est \emph{pas} au plus dénombrable. Expliquer en une phrase
        pourquoi, à travers la question 17, c'est exactement le théorème
        de Cantor (\cref{exo:b1:logic:11}) pour $E = \N$.
  \item Admettre --- comme on l'a vu au lycée, et comme on l'établira
        rigoureusement au \cref{ch:b1:reals} --- que tout $x \in
        \intco01$ possède un unique développement décimal \emph{propre}
        $x = 0.d_1 d_2 d_3\dots$ (c'est-à-dire ne se terminant pas par
        une infinité de $9$). Étant donnée une suite $(x_n)_{n \geq 1}$
        d'éléments de $\intco01$, construire $x \in \intco01$ tel que
        $x \neq x_n$ pour tout $n$~: choisir sa $n$-ième décimale égale
        à $5$ si la $n$-ième décimale de $x_n$ diffère de $5$, et à $6$
        sinon. Justifier soigneusement que $x$ est propre et évite tous
        les $x_n$, et en conclure que $\intco01$ n'est pas au plus
        dénombrable.
  \item En déduire que $\R$ n'est pas dénombrable, et que l'\omterm{def:b1:logic:sets}{ensemble}
        $\R \setminus \Q$ des nombres irrationnels ne l'est pas non
        plus. En quel sens précis «~la plupart~» des nombres réels
        sont-ils irrationnels~?
\end{enumerate}

\medskip
\noindent\textbf{Partie V --- Le théorème de Cantor de 1874~: les
\omterm{pb:b1:logic:1}{nombres transcendants} existent.} Un nombre réel $x$ est
\emph{algébrique}\index{nombre algébrique} lorsque $P(x) = 0$ pour un
certain polynôme non nul $P$ à coefficients entiers, et
\emph{transcendant}\index{nombre transcendant} sinon. Pour cette partie,
on admet --- c'est démontré au \cref{ch:b1:poly} --- qu'un polynôme non
nul de degré $n$ a au plus $n$ racines réelles.
\begin{enumerate}[resume]
  \item Montrer que tout nombre rationnel est algébrique, et trouver des
        polynômes explicites à coefficients entiers annulant $\sqrt 2$
        et $\sqrt 2 + \sqrt 3$.
  \item Pour $n \in \N$ fixé, montrer que l'\omterm{def:b1:logic:sets}{ensemble} des polynômes de
        degré au plus $n$ à coefficients entiers est dénombrable.
        (L'injecter dans $\Z^{n+1}$ et raisonner par récurrence sur $n$
        avec la question 13.)
  \item En déduire que l'\omterm{def:b1:logic:sets}{ensemble} de \emph{tous} les polynômes à
        coefficients entiers est dénombrable.
  \item Démontrer le \emph{théorème de Cantor sur les \omterm{pb:b1:logic:1}{nombres
        algébriques}}~: l'\omterm{def:b1:logic:sets}{ensemble} $\mathcal{A}$ des nombres réels
        algébriques est dénombrable.
  \item Conclure~: il existe des nombres réels \omterm{pb:b1:logic:1}{transcendants}, et
        l'\omterm{def:b1:logic:sets}{ensemble} des \omterm{pb:b1:logic:1}{nombres transcendants} n'est pas dénombrable.
        Faire ensuite le bilan de tout le problème en quelques phrases~:
        la chaîne $\N \approx \Z \approx \Q \approx \mathcal{A}$, le
        saut strict vers $\R \approx$ (essentiellement)
        $\mathcal{P}(\N)$, l'endroit où chaque outil
        (Cantor--Schröder--Bernstein, réunions dénombrables, la
        diagonale) a été décisif --- et la force philosophique d'une
        démonstration établissant l'existence d'une infinité non
        dénombrable de \omterm{pb:b1:logic:1}{nombres transcendants} sans en nommer un seul.
        (Démontrer qu'un nombre \emph{précis} comme $\pi$ est
        \omterm{pb:b1:logic:1}{transcendant} est bien plus difficile et dépasse ce volume.)
\end{enumerate}
\end{problem}
